\documentclass[10pt,a4paper]{article}
\usepackage[margin=18mm]{geometry}
\usepackage{amsmath,amssymb,booktabs,array,xcolor}
\usepackage[T1]{fontenc}
\setlength{\parindent}{0pt}
\setlength{\parskip}{4pt}
\newcommand{\GiB}{\,\mathrm{GiB}}
\newcommand{\MiB}{\,\mathrm{MiB}}

\title{Compact Resource Model for the Sparse Scenario QUBO}
\author{}
\date{29 August 2026}

\begin{document}
\maketitle
\vspace{-1.5em}

\section*{Definitions and assumptions}

\begin{tabular}{@{}>{$}l<{$} p{12.5cm}@{}}
A & number of assets; \\
z & binary states (z-bins) per asset; \\
k & directed top-$k$ neighbor nominations per asset; \\
m & reciprocal unordered nominations per asset; \\
u=k-m & unique undirected union pairs per asset, with $k/2\le u\le k$; \\
p & fraction of compatibility coefficients that are nonzero (approximately one); \\
r & simultaneous simulated-bifurcation runs; \\
b & scenarios packed into one block-diagonal solver batch; \\
\eta & usable-capacity safety factor (recommended: $\eta=0.75$).
\end{tabular}

The formulas assume the current undirected union graph, FP32 solver values,
32-bit CSR indices, FP64 compact CPU QUBO coefficients, and fixed $k$ as $A$
grows.  One pair is retained when either endpoint nominates the other; mutual
nominations are stored once.

\section*{Variables, interactions, and density}

The binary-variable count is
\begin{equation}
  N = Az.
\end{equation}

Upper-triangle QUBO interactions per asset are approximately
\begin{equation}
  e(z,k,u,p)
  = \underbrace{\frac{z(z-1)}{2}}_{\text{one-hot}}
  + \underbrace{puz^2}_{\text{compatibility}}.
  \label{eq:edges}
\end{equation}
Thus the upper interaction count is $E\simeq Ae$.  The symmetric solver CSR
matrix contains approximately $2Ae$ nonzeros.  QUBO density is
\begin{equation}
  \rho
  = \frac{E}{\binom{Az}{2}}
  \simeq \frac{2e}{Az^2}
  = O(A^{-1}).
\end{equation}
The graph therefore grows linearly in assets even though its variable-space
density decreases.

\section*{GPU VRAM model}

For FP32 values and 32-bit CSR columns, every upper interaction becomes two
symmetric entries, each containing a 4-byte value and 4-byte column index:
\begin{equation}
  M_{\mathrm{CSR}}/A \simeq 2e(4+4)=16e
  \quad\text{bytes per asset}.
\end{equation}

The measured FP32 solver model, including row pointers, fields, state tensors,
mask, force output, and sparse workspace, is
\begin{equation}
  g(z,e,r) \equiv M_{\mathrm{GPU}}/A
  \simeq 16e + 12z + 25zr
  \quad\text{bytes per asset}.
  \label{eq:gpu-per-asset}
\end{equation}
For a scenario batch of size $b$,
\begin{equation}
  M_{\mathrm{GPU}}
  \simeq M_{\mathrm{GPU,fixed}} + M_{\mathrm{corr}} + bA\,g(z,e,r),
\end{equation}
where $M_{\mathrm{corr}}$ is bounded approximately by
\texttt{--correlation-block-mib}.  The safe asset capacity is
\begin{equation}
  A_{\max,\mathrm{GPU}}
  \simeq
  \frac{\eta\,[V-M_{\mathrm{GPU,fixed}}-M_{\mathrm{corr}}]}
       {b\,g(z,e,r)}.
  \label{eq:gpu-capacity}
\end{equation}
Here $V$ is physical VRAM in bytes.  The number of solver steps does not enter
the memory formula because state arrays are reused at every step.

\section*{System RAM model}

The retained compact QUBO requires approximately
\begin{equation}
  M_{\mathrm{compact}}/A \simeq 16e+8z
  \quad\text{bytes per asset},
\end{equation}
from two 32-bit endpoints and one FP64 coefficient per upper interaction, plus
the FP64 linear vector.  Peak packing RAM is much larger because Ising,
symmetric COO, concatenation, and CSR arrays coexist transiently.  The measured
empirical model is
\begin{align}
  h(e) &\equiv M_{\mathrm{RAM,incremental}}/A
        \simeq \beta e, \\
  \beta &\simeq 130\text{--}135
        \quad\text{bytes per upper interaction}, \\
  M_{\mathrm{RAM}} &\simeq M_{\mathrm{RAM,fixed}}+bA\,h(e).
  \label{eq:ram-model}
\end{align}
The corresponding safe capacity is
\begin{equation}
  A_{\max,\mathrm{RAM}}
  \simeq
  \frac{\eta\,[H-M_{\mathrm{RAM,fixed}}]}
       {b\,\beta e},
  \label{eq:ram-capacity}
\end{equation}
where $H$ is RAM made available to the process.  This coefficient is empirical
and should be recalibrated after changes to CSR packing.

\section*{Calibration at current defaults}

For $A=10{,}000$, $z=21$, $k=5$, $u\simeq4.16$, $p\simeq1$,
$r=16$, and $b=1$:
\begin{align}
  N &= 210{,}000, \\
  E &= 20{,}433{,}166, \\
  e &= 2{,}043.3166, \\
  \rho &\simeq 9.27\times10^{-4}\;(0.0927\%), \\
  M_{\mathrm{compact}} &= 313.4\MiB, \\
  M_{\mathrm{GPU,peak}} &= 402.8\MiB, \\
  M_{\mathrm{GPU,reserved}} &= 420\MiB, \\
  M_{\mathrm{process,peak}} &= 3.22\GiB.
\end{align}

The measured GPU allocation is about $42.2$ kB per asset.  Incremental peak
process RAM above imported-library baseline is about $265$ KiB per asset.
Therefore:
\begin{align}
  \text{raw VRAM capacity} &\simeq 25{,}400
       &&\text{assets per available GiB}, \\
  \text{safe VRAM capacity} &\simeq 19{,}000
       &&\text{assets per available GiB}, \\
  \text{raw RAM capacity} &\simeq 4{,}000
       &&\text{assets per available GiB}, \\
  \text{safe RAM capacity} &\simeq 3{,}000
       &&\text{assets per available GiB}.
\end{align}

\begin{center}
\begin{tabular}{@{}r rr@{}}
\toprule
Usable resource & VRAM-limited assets & RAM-limited assets \\
\midrule
$1\GiB$   & 19,000  & 3,000 \\
$8\GiB$   & 152,000 & 24,000 \\
$16\GiB$  & 304,000 & 48,000 \\
$32\GiB$  & 608,000 & 96,000 \\
$64\GiB$  & ---     & 192,000 \\
$128\GiB$ & ---     & 384,000 \\
\bottomrule
\end{tabular}
\end{center}

\section*{Measured effect of simultaneous runs}

\begin{center}
\begin{tabular}{@{}r rr@{}}
\toprule
Runs $r$ & Peak allocated & Peak reserved \\
\midrule
1  & $327.5\MiB$ & $342\MiB$ \\
4  & $353.6\MiB$ & $358\MiB$ \\
16 & $402.8\MiB$ & $420\MiB$ \\
32 & $485.6\MiB$ & $512\MiB$ \\
\bottomrule
\end{tabular}
\end{center}

Total runs may exceed the simultaneous width without increasing peak state
memory by setting \texttt{--run-batch-size}.  For example, 32 total runs with a
run batch of 16 retain approximately the 16-run state-memory requirement but
require two solver passes.

\section*{Scaling rules}

\begin{center}
\begin{tabular}{@{}ll@{}}
\toprule
Parameter change & Leading memory effect \\
\midrule
$A$ assets & linear \\
$k$ neighbors & approximately linear through $u$ \\
$z$ bins & approximately quadratic through $uz^2$ \\
$r$ simultaneous runs & linear state-memory term $25zr$ \\
$b$ packed scenarios & approximately linear in graph and state memory \\
solver steps & no material peak-memory change \\
FP64 solver & larger value/state terms; indices unchanged \\
\bottomrule
\end{tabular}
\end{center}

In particular, changing $z$ from 21 to 41 changes the dominant compatibility
term by approximately
\begin{equation}
  \left(\frac{41}{21}\right)^2 \simeq 3.81.
\end{equation}

\section*{Implementation ceiling and caveats}

The current union graph produces approximately $2e\simeq4{,}087$ symmetric CSR
nonzeros per asset.  Signed 32-bit CSR indexing therefore implies an approximate
ceiling
\begin{equation}
  A_{\max,\mathrm{index}}
  \simeq \frac{2^{31}-1}{2e}
  \simeq 525{,}000\ \text{assets}.
\end{equation}
This may bind before VRAM on large-memory accelerators.  Actual deployments
must also reserve memory for CUDA context, display/other processes, correlation
blocks, the operating system, and allocator fragmentation.  Runtime and CPU
RAM may become impractical well before a VRAM-only limit is reached.

\end{document}
